\documentclass[11pt,a4paper]{article}
\usepackage[T1]{fontenc}
\usepackage{lmodern,amsmath,amssymb,booktabs,microtype}
\usepackage[margin=24mm]{geometry}
\usepackage[colorlinks=true,urlcolor=blue,linkcolor=blue]{hyperref}
\newcommand{\EE}{\mathbb E}
\newcommand{\core}{\mathcal E}
\DeclareMathOperator{\tr}{tr}
\DeclareMathOperator{\Rem}{Rem}
\title{Charge removal opens product-forbidden rank channels\\An exact correction to the memory-filter ansatz}
\author{Alan Day \and Chengheng\thanks{AI research assistant (OpenAI).}}
\date{Local research note --- 6 November 2026}
\begin{document}\maketitle
\begin{abstract}
The collinearity of four quadratic maps in the previously tested
$N=5,p=2,m=3$ system does not extend to general charge sectors.
At half filling, particle-hole transposition imposes an exact parity
selection rule on symmetrized products. Charge removal can populate
the forbidden channels. We give an analytic family for the force norm
$A$ and an independently verified counterexample for the actual
quadratic contraction of the quartic memory core. A two-charge
amplitude representation and an exact norm identity for charge
removal supply the corrected starting point for the spectral limit.
\end{abstract}

\section{A selection rule for the multiplication map}
Use standard CAR and the full overlapping complex Gaussian $p$-body
Hamiltonian, $\EE J_{IJ}J_{KL}=w\delta_{IL}\delta_{JK}$, with
$\hbar=1$. On charge $m$, $\tr_m$ denotes normalized trace and $P_r$
the neutral $U(N)$ rank projector. All statements through Section 4
are finite-dimensional identities, with no double-scaling limit.

For $N=2m$, let $\star$ be the signed complement map on
$\wedge^m\mathbb C^N$, and define the linear anti-automorphism
\begin{equation}
 \tau(X)=\star X^{\mathsf T}\star^{-1},\qquad
 \tau(XY)=\tau(Y)\tau(X),\qquad \tau^2=1.
\end{equation}
It sends $E_{ij}=c_i^\dagger c_j$ to $\delta_{ij}-E_{ij}$.
Consequently it commutes with the adjoint $U(N)$ action.
The neutral operator representation is multiplicity-free in rank.
On a disjoint, traceless rank-$r$ highest-weight monomial, reversing
and normal ordering gives $(-1)^r$. Thus
\begin{equation}
 \tau P_r=(-1)^r P_r.
 \label{eq:parity}
\end{equation}
For arbitrary pure-rank $X_a,Y_b$, it follows that
\begin{equation}
 P_r\{X_a,Y_b\}=0\quad\hbox{if }a+b-r\text{ is odd}.
 \label{eq:zero}
\end{equation}
This is a statement about the multiplication map, not the absence of
the irrep from the tensor product. In particular, every odd-rank
component of $H_a^2$ vanishes. An equivariant quadratic map need not
satisfy the same selection rule.

\newpage
\section{A family of pure-rank counterexamples}
Let $n_i=c_i^\dagger c_i$, $1\le p<m=N/2$, and choose the deterministic
$p$-body input
\begin{equation}
 H=\prod_{a=1}^p d_a,\qquad
 d_a=n_{2a-1}-n_{2a},\quad
 t_a=n_{2a-1}+n_{2a},\quad b_a=d_a^2.
\end{equation}
Every partial contraction of its $p$-body coefficient tensor vanishes:
the two choices in the omitted pair cancel. Hence this input is pure
rank $p$ wherever its embedding is nonzero. Also
$\tau H=(-1)^pH$ and $H^2=\prod_a b_a$ is even.

Direct removal of an occupied mode gives
\begin{equation}
 A=\sum_i[H,c_i]^\dagger[H,c_i]
   =\sum_{a=1}^p t_a\prod_{b\ne a}b_b.
\end{equation}
Under complementation, $t_a\mapsto2-t_a$ and $b_a\mapsto b_a$.
Its odd part is therefore
\begin{equation}
 A_{\rm odd}:=\frac{A-\tau A}{2}
       =\sum_{a=1}^p(t_a-1)\prod_{b\ne a}b_b\ne0.
 \label{eq:odd}
\end{equation}
A configuration with one double-occupied pair and all other selected
pairs singly occupied has value $+1$; an empty exceptional pair gives
$-1$. The $N-2p$ remaining modes can accommodate the required charge
because $m>p$. Every other configuration gives zero.

Counting these configurations yields exact normalized HS norms:
\begin{align}
 \|H^2\|_m^2
 &=\frac{2^p\binom{N-2p}{m-p}}{\binom Nm},\\
 \|A_{\rm odd}\|_m^2
 &=\frac{p\,2^p\binom{N-2p}{m-p-1}}{\binom Nm},\\
 \frac{\|A_{\rm odd}\|_m^2}{\|H^2\|_m^2}
 &=p\,\frac{m-p}{m-p+1}.
 \label{eq:ratio}
\end{align}
Thus no scalar coefficients multiplying the channels $P_rH_p^2$
can reproduce $A$ for this family. It exists along half-filled
double-scaled sequences as well, but these are deterministic inputs;
Eq.~\eqref{eq:ratio} is not an annealed memory-strength calculation.

The checker exhausts all charge-$m$ configurations in 15 systems,
$N=4,6,8,10,12$ and every $1\le p<m$. It verifies the removal formula,
both norm counts, and the rank eigenvalue
$\mathcal C H=2p(N+1-p)H$ using the diagonal Casimir action.

\newpage
\section{The actual quadratic memory contraction}
The target is still the previously defined core
\begin{equation}
 \core=\sum_i(\operatorname{ad}_H^2c_i)^\dagger
                   \operatorname{ad}_H^2c_i
                  +\tfrac14[H,[H,A]],\qquad
 Q_{\core}=\tfrac12\Delta\core.
\end{equation}
For this homogeneous quartic polynomial, ordinary Gaussian chaos
gives $W_2\core=Q_{\core}-2\EE\core$. The ensemble mean is scalar in
each charge sector, so $P_1W_2\core=P_1Q_{\core}$.

Take $N=6,p=2,m=3,w=1$ and
\begin{equation}
 h=(n_1-n_2)(n_3-n_4)\big|_3,\qquad
 k=(n_1-n_2)(n_3-n_4)\big|_2,\qquad
 D=n_1+n_2+n_3+n_4-2I.
\end{equation}
Both inputs have Casimir eigenvalue $20$, hence rank two. Exact
projection gives
\begin{equation}
 \boxed{P_1h^2=0,\qquad P_1A=D,\qquad P_1Q_{\core}=99D.}
 \label{eq:witness}
\end{equation}
Here $D$ has eigenvalues $+1$ on four states, $-1$ on four states,
and zero on twelve states. In the normalized charge-three trace,
\begin{equation}
 \|P_1A\|_3^2=\frac25,\qquad
 \|P_1Q_{\core}\|_3^2=\frac{19602}{5}.
\end{equation}
The old $J_hF$ direction cannot repair this zero within the $(2,2;1)$
block: on a pure input rank, the equivariant linear $F$ is a scalar
multiple of $h$, so $P_1J_hF=0$ as well.

The calculation uses the existing explicit $\Phi_\pm,\Psi,\Rem$
single-contraction formula. Independently, full quartic polarization
of the commutator-defined core over all 225 real Gaussian coordinates
reproduces $Q_{\core}$ exactly. The raw-word identity is checked too.
Four full matrix Casimir projectors verify
Eq.~\eqref{eq:parity}, not only its diagonal restriction.

This disproves the general scalar-filter ansatz. It does not disprove
the previous $N=5$ calculation or its exact Gram reconstruction.
The present diagonal realization commutes with its projected $D$;
it is not a new noncommuting dynamical witness. Neither a general
multiplicity rank nor the double-scaled stochastic weight follows
from this one realization.

\newpage
\section{The corrected two-charge interface}
Write $h_a=H_m^{(a)}$, $k_a=H_{m-1}^{(a)}$ using the same coupling
irrep. Define symmetric products as $X\odot Y=\{X,Y\}$ for distinct
input ranks, and $X\odot X=X^2$ for equal ranks. Keep both maps
\begin{equation}
 Z^+_{ab;r}=P_r(h_a\odot h_b),\qquad
 Z^-_{ab;r}=P_r\Rem(k_a\odot k_b).
\end{equation}
Then the exact old identity $A=\Rem(k^2)+(2p-m)h^2$ becomes
\begin{equation}
 A_{ab;r}=Z^-_{ab;r}+(2p-m)Z^+_{ab;r}.
\end{equation}
No division by $Z^+$ is allowed. In reduced-tensor notation, choose
primitive normal-ordered rank tensors $T_a^{(m)}$ with consistent
normalization across charges, so
$h_a=e_a(m)T_a^{(m)}$, $e_a(m)=\binom{m-a}{p-a}$.
If $K^{(m)}_{ab;r,\rho}$ denotes their symmetrized product coefficient
in a common intertwiner basis, the amplitude of $A$ is
\begin{align}
 \mathsf A_{ab;r,\rho}
 ={}&(m-r)e_a(m-1)e_b(m-1)K^{(m-1)}_{ab;r,\rho}\notag\\
 &+(2p-m)e_a(m)e_b(m)K^{(m)}_{ab;r,\rho}.
 \label{eq:twocharge}
\end{align}
This uses $\Rem(T_r^{(m-1)})=(m-r)T_r^{(m)}$. The two reduced
coefficients include their relative signs and multiplicity label
$\rho$; the recoupling framework already exists in
Kota--Vyas~\cite{KV}. Their dilute limit is not the present limit.

There is also an exact norm transfer for every rank-$r$ lower-sector
operator $X_r$ in the common rank range:
\begin{equation}
 \|\Rem X_r\|_m^2=\ell_r^2\|X_r\|_{m-1}^2,\qquad
 \ell_r^2=\frac{m(m-r)(N-m+1-r)}{N-m+1}\le m^2.
 \label{eq:lift}
\end{equation}
To prove it, a disjoint rank-$r$ monomial has norm proportional to
$\binom{N-2r}{m-r}/\binom Nm$. Taking the ratio between adjacent
charges and multiplying by $(m-r)^2$ gives Eq.~\eqref{eq:lift}.
Equivariance and irreducibility extend it to the full rank space.
Eight direct CAR checks confirm it. Components outside the common
rank range are killed by the equivariant removal map.

For $a<b$, the previous positive ordered-product law now gives
\begin{equation}
 \EE\|Z^-_{ab;r}\|_m^2
 \le4\ell_r^2 C_a^-C_b^-\pi^{(m-1)}_{ab}(r),
 \qquad C_a^-I=\EE k_a^2.
\end{equation}
For $a=b$ after $W_2$ centering, replace $4C_a^-C_b^-$ by
$2(C_a^-)^2$. These are bounds; the same-coupling cross term
$\EE\langle Z^+,Z^-\rangle_m$ remains essential. The next calculation
is therefore the coherent two-charge amplitudes of $Q_{\core}$,
including $\Psi$ insertions, before taking rank-deficit limits.

\begin{thebibliography}{9}
\bibitem{KV} V. K. B. Kota and M. Vyas,
\emph{Random Matrix Theory for Transition Strength Densities in Finite
Quantum Systems}, Sec.~3, Eqs.~(34)--(35),
\href{https://arxiv.org/abs/1504.00759}{arXiv:1504.00759}.
\bibitem{Kota} V. K. B. Kota,
\emph{Bivariate moments of the two-point correlation function for
embedded Gaussian unitary ensemble with $k$-body interactions},
Sec.~IV, \href{https://arxiv.org/abs/2208.11312}{arXiv:2208.11312v2}.
\end{thebibliography}
\end{document}
